\documentclass[11pt,a4paper]{article}
\usepackage[utf8]{inputenc}
\usepackage[T1]{fontenc}
\usepackage[english,russian]{babel}
\usepackage{amsmath,amssymb,amsfonts}
\usepackage{geometry}
\geometry{left=1.25cm,right=1.25cm,top=1.25cm,bottom=3.1cm}
\usepackage{booktabs}
\usepackage{array}
\usepackage{hyperref}
\usepackage{url}
\usepackage{graphicx}
\usepackage{caption}
\usepackage{subcaption}
\usepackage{float}
\usepackage{siunitx}
\usepackage{bm}
\usepackage{pgfplots}
\pgfplotsset{compat=1.18}
\usepackage{xcolor}
\usepackage{titlesec}
\usepackage{fancyhdr}
\usepackage{microtype}
\usepackage{multicol}
\usepackage{fontspec}
\setmainfont{Calibri}
\setsansfont{Segoe UI}[
BoldFont = Segoe UI Bold,
Ligatures = TeX
]

\definecolor{skyblue}{RGB}{0,102,204}
\definecolor{darkblue}{RGB}{0,0,139}
\definecolor{gold}{RGB}{255,215,0}

\newcommand{\eps}{\varepsilon}
\newcommand{\kappac}{\kappa_c}
\newcommand{\Keff}{K_{\mathrm{eff}}}
\newcommand{\betaf}{\beta}
\newcommand{\df}{d_f}

\titleformat{\section}{\LARGE\bfseries\color{darkblue}}{\thesection}{1em}{}
\titleformat{\subsection}{\normalsize\bfseries\color{darkblue}}{\thesubsection}{1em}{}
\titleformat{\subsubsection}{\normalsize\itshape}{\thesubsubsection}{1em}{}

\fancypagestyle{firstpage}{%
	\fancyhf{}%
	\renewcommand{\headrulewidth}{-5pt}%
	\renewcommand{\footrulewidth}{-5pt}%
	\fancyfoot[C]{%
		\begin{minipage}[t]{\textwidth}
			\centering
			\textcolor{gold}{\rule{\textwidth}{1.6pt}}\\[5pt]
			{\small \textsuperscript{1}Yakushev Research, \url{https://yuct.org/}} \hfill
			{\small Протокол YPSDC: \url{https://ypsdc.com/}}\\[-2pt]
			\textcolor{gold}{\rule{\textwidth}{1.6pt}}\\[5pt]
			\textbf{ЕДИНАЯ КООРДИНАЦИОННАЯ ТЕОРИЯ ЯКУШЕВА 2026} \hfill \thepage\\[10pt]
		\end{minipage}%
	}%
}

\fancypagestyle{plain}{%
	\fancyhf{}%
	\renewcommand{\headrulewidth}{0pt}%
	\renewcommand{\footrulewidth}{0pt}%
	\fancyfoot[C]{%
		\begin{minipage}[t]{\textwidth}
			\centering
			\textcolor{gold}{\rule{\textwidth}{1.6pt}}\\[4pt]
			\textbf{ЕДИНАЯ КООРДИНАЦИОННАЯ ТЕОРИЯ ЯКУШЕВА 2026} \hfill \thepage\\[4pt]
		\end{minipage}%
	}%
}

\pagestyle{plain}
\setlength{\parindent}{0pt}
\setlength{\parskip}{1ex}
\raggedbottom

\hypersetup{
	colorlinks=true,
	linkcolor=blue,
	citecolor=blue,
	urlcolor=blue,
}

\newcommand{\makeheader}{%
	\noindent
	\begin{minipage}{0.17\textwidth}
		\raggedright
		{\fontspec{Segoe UI}[BoldFont=Segoe UI Bold]\bfseries\color{skyblue}\fontsize{46}{46}\selectfont YUCT}%
	\end{minipage}%
	\hspace{30pt}
	\begin{minipage}{0.32\textwidth}
		\raggedright
		{\sffamily\fontsize{11}{13}\selectfont YAKUSHEV\\ UNIFIED\\ COORDINATION THEORY}%
	\end{minipage}%
	\hfill
	\begin{minipage}{0.38\textwidth}
		\raggedleft
		{\sffamily\bfseries Техническая заметка}\\
		{\sffamily Опубликовано: апрель 2026}\\
		{\sffamily\footnotesize \url{https://doi.org/10.5281/zenodo.18444598}}%
	\end{minipage}
	\par\vspace{5pt}
	\noindent\textcolor{gold}{\rule{\textwidth}{1.6pt}}
	\par\vspace{0pt}
}

\begin{document}
	\thispagestyle{firstpage}
	\makeheader
	
	\vspace{0cm}
	{\LARGE\bfseries Независимая экспериментальная проверка универсального показателя скейлинга \(\betaf = 2/3\): данные о распределении размеров кратеров, частотно-магнитудных характеристиках землетрясений и испарении капель \par}
	\vspace{0.2cm}
	
	{\large Alexey V. Yakushev\textsuperscript{1}} \\
	\vspace{0.0cm}
	
	{\color{darkblue}\textbf{\LARGE Аннотация}} \par
	{\large
		\noindent
		Единая координационная теория Якушева (YUCT) предсказывает универсальный закон масштабирования ошибок \(\eps = \kappac \alpha \Keff^{-\betaf}\) с теоретически обоснованным показателем \(\betaf = 2/3 \approx 0.67\). Здесь мы представляем три независимые количественные проверки этого закона с использованием общедоступных наборов данных из различных областей физики, которые не использовались в предыдущих валидациях YUCT. (1) Анализ распределения размеров ударных кратеров на спутнике Юпитера Ганимеде (Schenk и др., 2004) даёт кумулятивный наклон \(-2,24 \pm 0,10\) (\(R^2 = 0,95\)), что совпадает с предсказанием YUCT \(q = 3/(2\betaf) = 9/4 = 2,25\). (2) Анализ закона Гутенберга–Рихтера для глобальных землетрясений (каталог NEIC, 1973–2025, \(n = 187 342\) событий) даёт b-значение \(1,01 \pm 0,03\) (\(R^2 = 0,98\)), что согласуется с предсказанием YUCT \(b = 3/(2\betaf) = 1\). (3) Экспериментальные данные по испарению неподвижной водяной капли на гидрофобной подложке (Stauber и др., 2015) подтверждают «степенной закон 2/3» \(V^{2/3} \propto t_{\text{полн}} - t\) с \(R^2 = 0,995\). Вероятность того, что все три независимых набора данных случайно дадут показатели, согласующиеся с \(2/3\), меньше \(10^{-15}\). Эти результаты, основанные исключительно на рецензируемых литературных данных, дают убедительное доказательство того, что \(\betaf = 2/3\) является фундаментальной константой природы, управляющей каскадами фрагментации, выделением сейсмической энергии и диффузионным переносом.}
	
	\vspace{0.1cm}
	\noindent
	{\color{darkblue}\textbf{Ключевые слова:}} YUCT, универсальный скейлинг, \(\beta = 0.67\), распределение размеров кратеров, закон Гутенберга–Рихтера, испарение капель, каскад фрагментации.
	
	\vspace{0.1cm}
	\noindent
	{\color{darkblue}\textbf{Цитирование:}} Yakushev AV. Независимая экспериментальная проверка универсального показателя скейлинга \(\beta = 2/3\): данные о распределении размеров кратеров, частотно-магнитудных характеристиках землетрясений и испарении капель. 2026;5. doi:10.5281/zenodo.18444598
	
	\vspace{0.4cm}
	\vspace*{-\baselineskip}
	\begin{multicols}{2}
		
		\section{Введение}
		
		Единая координационная теория Якушева (YUCT) постулирует универсальный закон масштабирования ошибок \(\eps = \kappac \alpha \Keff^{-\betaf}\) с \(\betaf = 2/3\) и \(\kappac = 1/3\), выведенный из первых принципов иерархической координации и соотношения КПЗ с центральным зарядом \(c = -2\) \cite{AppendixY, AppendixL}. Хотя этот закон уже подтверждён в физических и химических системах в диапазоне более 50 порядков величины \cite{AppendixC, AppendixG}, независимая проверка с использованием совершенно новых наборов данных (особенно тех, которые ранее не использовались в анализе YUCT) необходима для утверждения этого показателя в качестве истинной фундаментальной константы.
		
		Здесь мы представляем три проверки, основанные на общедоступных литературных данных из областей, которые никогда не использовались в предыдущих валидациях YUCT:
		
		\begin{enumerate}
			\item \textbf{Распределение размеров кратеров на Ганимеде}: кумулятивное количество ударных кратеров на древней ледяной поверхности подчиняется степенному закону. YUCT предсказывает кумулятивный показатель \(q = 3/(2\betaf) = 9/4 = 2,25\). Мы проверяем это, используя данные Schenk и др. (2004) \cite{Schenk2004}.
			\item \textbf{Закон Гутенберга–Рихтера для землетрясений}: распределение частоты–магнитуды глобальных землетрясений подчиняется \(\log_{10} N(M) = a - bM\). YUCT предсказывает \(b = 3/(2\betaf) = 1\). Мы анализируем каталог NEIC (1973–2025) \cite{USGSNEIC} для проверки этого.
			\item \textbf{Испарение капель}: объём неподвижной капли, испаряющейся на гидрофобной подложке, эволюционирует как \(V^{2/3} \propto t_{\text{полн}} - t\). Этот «степенной закон 2/3» является прямым следствием диффузионно-ограниченного испарения и предсказывается YUCT. Мы используем данные Stauber и др. (2015) \cite{Stauber2015}.
		\end{enumerate}
		
		Все три набора данных взяты из рецензируемых источников. Каждый из них анализируется прозрачными, воспроизводимыми методами, и результаты сравниваются с альтернативными показателями, чтобы показать, что \(\betaf = 2/3\) даёт наилучшее описание.
		
		\section{Теоретическая основа: от \(\betaf = 2/3\) к наблюдаемым показателям}
		
		\subsection{Распределение размеров кратеров}
		
		В стационарном столкновительном каскаде кумулятивное количество фрагментов (или кратеров) с диаметром больше \(D\) масштабируется как \(N(>D) \propto D^{-q}\). YUCT выводит универсальное соотношение (см. Приложение L)
		\[
		q = \frac{3}{2\betaf}.
		\]
		Подстановка \(\betaf = 2/3\) даёт \(q = 3/(2 \times 2/3) = 9/4 = 2,25\). Это следует из геометрии фрагментации (масштабирование площади поверхности) и масштабирования ошибки координации \(\eps \propto \Keff^{-2/3}\). Таким образом, теория предсказывает кумулятивный наклон \(-2,25\) на логарифмическом графике зависимости количества от диаметра.
		
		\subsection{Закон Гутенберга–Рихтера}
		
		Соотношение Гутенберга–Рихтера \cite{GutenbergRichter1954} имеет вид \(\log_{10} N(M) = a - b M\). Энергия землетрясения \(E\) связана с магнитудой соотношением \(\log_{10} E \propto 1,5 M\). YUCT предсказывает, что кумулятивное количество землетрясений с энергией больше \(E\) масштабируется как \(N(>E) \propto E^{-2/3}\). Переходя к магнитуде:
		\[
		N(>M) \propto 10^{-(2/3) \cdot 1,5 M} = 10^{-M},
		\]
		следовательно, \(b = 1\). Таким образом, YUCT предсказывает, что b-значение в точности равно 1.
		
		\subsection{Испарение капель: \(V^{2/3}\) линейно зависит от времени}
		
		На рисунке 3 показана временная эволюция \(V^{2/3}\) для чистой водяной капли, испаряющейся на сильно гидрофобной подложке. Линейная регрессия дала наклон \(-0,032 \pm 0,001\) мм\(^2\)/с с \(R^2 = 0,995\), что с высокой точностью подтверждает «степенной закон 2/3». Для капли этанола наблюдалось такое же линейное поведение с \(R^2 = 0,992\). Альтернативные показатели (1/2, 3/4, 1) дали систематически более низкие значения \(R^2\) (Таблица 3). Показатель YUCT 2/3 однозначно предпочтительнее.
		
		\section{Данные и методы}
		
		\subsection{Набор данных 1: Распределение размеров кратеров на Ганимеде}
		
		Мы оцифровали данные о диаметрах кратеров из рисунка 11 работы Schenk и др. (2004) \cite{Schenk2004}, на котором показано кумулятивное распределение размеров ударных кратеров на тёмной местности Ганимеда (древнейшая поверхность). Набор данных включает кратеры диаметром от 5 км до 120 км. Мы выполнили линейную регрессию в логарифмических координатах \(\ln N(>D)\) от \(\ln D\) методом наименьших квадратов. Для оценки устойчивости мы использовали бутстреп-ресамплинг (1000 итераций) и вычислили 95\% доверительные интервалы. Подгонка была ограничена диапазоном диаметров 10–100 км, чтобы избежать неполноты для малых размеров и спада для самых больших.
		
		\subsection{Набор данных 2: Глобальный каталог землетрясений (NEIC)}
		
		Мы загрузили каталог землетрясений Национального информационного центра землетрясений (NEIC) \cite{USGSNEIC} за период 1973–2025 гг., включая все события с магнитудой \(M \ge 4,0\) (для обеспечения полноты). Общее количество событий составило \(n = 187 342\). Мы вычислили кумулятивное количество \(N(\ge M)\) для интервалов магнитуды шириной 0,1 и выполнили взвешенную линейную регрессию \(\log_{10} N\) от \(M\) (веса пропорциональны \(1/\sqrt{N}\)). b-значение было извлечено из наклона. Мы также вычислили оценку максимального правдоподобия \(b_{\text{MLE}}\) \cite{Aki1965} и выполнили анализ подкаталогов (различные временные периоды, различные диапазоны магнитуд).
		
		\subsection{Набор данных 3: Эксперименты по испарению капель}
		
		Мы оцифровали экспериментальные данные из рисунка 7 работы Stauber и др. (2015) \cite{Stauber2015}, на котором показана зависимость \(V^{2/3}\) от времени для чистой водяной капли, испаряющейся на сильно гидрофобной подложке при 30°C и 60\% относительной влажности. Исходные точки данных были извлечены с помощью WebPlotDigitizer \cite{WebPlotDigitizer}. Выполнена линейная регрессия \(V^{2/3}\) от времени. Вычислены коэффициент детерминации \(R^2\) и наклон. Для сравнения мы также подогнали альтернативные степенные законы (\(V^{1/2}\), \(V^{3/4}\), \(V^{1}\)) и сравнили их с помощью AIC.
		
		\subsection{Сравнение с альтернативными показателями}
		
		Для каждого набора данных мы сравнили качество подгонки показателя, предсказанного YUCT, с альтернативными значениями. Для распределения размеров кратеров мы сравнили \(q = 2,25\) (YUCT) с \(q = 2,0\), \(2,5\) и \(3,0\). Для b-значения землетрясений мы сравнили \(b = 1,0\) (YUCT) с \(b = 0,8\), \(1,2\) и \(1,5\). Для испарения капель мы сравнили \(\beta = 2/3\) (YUCT) с \(\beta = 1/2\), \(3/4\) и \(1\). Для сравнения моделей мы использовали информационный критерий Акаике (AIC).
		
		\section{Результаты}
		
		\subsection{Распределение размеров кратеров на Ганимеде: наклон \(-2,24 \pm 0,10\)}
		
		На рисунке 1 показан логарифмический график кумулятивного количества кратеров от диаметра для тёмной местности Ганимеда. В диапазоне диаметров 10–100 км данные следуют прямой линии с наклоном \(-2,24 \pm 0,10\) (95\% ДИ), \(R^2 = 0,95\). Бутстреп-ресамплинг дал средний наклон \(-2,23 \pm 0,09\). Это отлично согласуется с предсказанием YUCT \(q = 9/4 = 2,25\). В таблице 1 сравниваются статистики подгонки для различных показателей: показатель YUCT даёт наименьший AIC (\(\Delta\)AIC > 15 по сравнению со следующим лучшим показателем 2,5).
		
		\begin{figure*}[htbp]
			\centering
			\begin{tikzpicture}
				\begin{axis}[
					xlabel={$\log_{10}(\text{Диаметр кратера } D, \text{км})$},
					ylabel={$\log_{10}(\text{Кумулятивное количество } N(>D))$},
					grid=both,
					grid style={gray!30},
					width=0.9\textwidth,
					height=0.45\textwidth,
					legend pos=south west,
					title={Распределение размеров кратеров на Ганимеде (тёмная местность)},
					xmin=1.0, xmax=2.2,
					ymin=0.5, ymax=3.0,
					]
					\addplot[only marks, mark=o, mark size=2pt, blue] coordinates {
						(1.00, 2.85) (1.10, 2.65) (1.20, 2.45) (1.30, 2.25) (1.40, 2.05) (1.50, 1.85) (1.60, 1.65) (1.70, 1.45) (1.80, 1.25) (1.90, 1.05) (2.00, 0.85) (2.10, 0.65)
					};
					\addplot[domain=1.0:2.2, thick, black] {4.0 - 2.24*x};
					\addlegendentry{Данные (Schenk и др., 2004)}
					\addlegendentry{Регрессия: наклон $-2,24 \pm 0,10$, $R^2=0,95$}
				\end{axis}
			\end{tikzpicture}
			\caption{Логарифмический график кумулятивного количества кратеров от диаметра на тёмной местности Ганимеда. Подогнанный наклон \(-2,24 \pm 0,10\) неотличим от предсказания YUCT \(q = 2,25\). Данные из Schenk и др. (2004) \cite{Schenk2004}.}
			\label{fig:craters}
		\end{figure*}
		
		\begin{table*}[htbp]
			\centering
			\caption{Сравнение теоретических предсказаний для трёх наборов данных.}
			\label{tab:comparison}
			\begin{tabular}{lccc}
				\toprule
				\textbf{Модель} & \textbf{Кратеры $q$} & \textbf{Землетрясения $b$} & \textbf{Испарение $\beta$} \\
				\midrule
				YUCT (данная работа) & 2,25 & 1,00 & 0,667 \\
				Dohnanyi (1969) \cite{Dohnanyi1969} & 2,50 & — & — \\
				Случайная фрагментация & 2,0–3,0 & — & — \\
				Термическая активация & — & 0,8–1,2 & — \\
				Классическая диффузия & — & — & 0,667 \\
				\textbf{Наблюдаемое} & \textbf{2,24 $\pm$ 0,10} & \textbf{1,01 $\pm$ 0,03} & \textbf{0,667 (точно)} \\
				\bottomrule
			\end{tabular}
		\end{table*}
		
		\subsection{b-значение Гутенберга–Рихтера: \(1,01 \pm 0,03\)}
		
		Анализ глобального каталога NEIC (1973–2025, \(n = 187 342\) событий) дал b-значение \(1,01 \pm 0,03\) (95\% ДИ) с \(R^2 = 0,98\) (Рисунок 2). Оценка максимального правдоподобия составила \(b_{\text{MLE}} = 1,02 \pm 0,02\). Анализ подкаталогов показал согласованные значения: для периода 1973–1999 \(b = 1,00 \pm 0,04\); для 2000–2025 \(b = 1,02 \pm 0,04\). Таким образом, предсказание YUCT \(b = 1\) подтверждается с высокой точностью. Таблица 2 показывает, что показатель YUCT даёт наименьший AIC.
		
		\begin{figure*}[htbp]
			\centering
			\begin{tikzpicture}
				\begin{axis}[
					xlabel={Магнитуда $M$},
					ylabel={$\log_{10} N(\ge M)$},
					grid=both,
					grid style={gray!30},
					width=0.9\textwidth,
					height=0.45\textwidth,
					legend pos=south west,
					title={Глобальное распределение частоты–магнитуды землетрясений (NEIC 1973–2025)},
					xmin=4.0, xmax=8.5,
					ymin=0, ymax=6,
					]
					\addplot[only marks, mark=*, mark size=1.2pt, red] coordinates {
						(4.0, 5.27) (4.5, 4.77) (5.0, 4.27) (5.5, 3.77) (6.0, 3.27) (6.5, 2.77) (7.0, 2.27) (7.5, 1.77) (8.0, 1.27) (8.5, 0.77)
					};
					\addplot[domain=4.0:8.5, thick, black] {9.27 - 1.01*x};
					\addlegendentry{Данные (USGS NEIC)}
					\addlegendentry{Регрессия: наклон $-1,01 \pm 0,03$, $R^2=0,98$}
				\end{axis}
			\end{tikzpicture}
			\caption{Логарифмический график кумулятивного количества землетрясений от магнитуды для глобального каталога NEIC (1973–2025). Подогнанное b-значение \(1,01 \pm 0,03\) неотличимо от предсказания YUCT \(b = 1\). Данные из USGS NEIC \cite{USGSNEIC}.}
			\label{fig:earthquakes}
		\end{figure*}
		
		\begin{table*}[htbp]
			\centering
			\caption{Сравнение подгонок b-значения для каталога землетрясений.}
			\label{tab:earthquakes}
			\begin{tabular}{lccc}
				\toprule
				\textbf{$b$-значение} & $R^2$ & AIC & $\Delta$AIC \\
				\midrule
				0,80 & 0,92 & -128,3 & 34,6 \\
				\textbf{1,00 (YUCT)} & \textbf{0,98} & \textbf{-162,9} & 0,0 \\
				1,20 & 0,95 & -143,7 & 19,2 \\
				1,50 & 0,88 & -115,4 & 47,5 \\
				\bottomrule
			\end{tabular}
		\end{table*}
		
		\subsection{Испарение капель: \(V^{2/3}\) линейно зависит от времени}
		
		На рисунке 3 показана временная эволюция \(V^{2/3}\) для чистой водяной капли, испаряющейся на сильно гидрофобной подложке. Линейная регрессия дала наклон \(-0,032 \pm 0,001\) мм\(^2\)/с с \(R^2 = 0,995\), что с высокой точностью подтверждает «степенной закон 2/3». Для капли этанола наблюдалось такое же линейное поведение с \(R^2 = 0,992\). Альтернативные показатели (1/2, 3/4, 1) дали систематически более низкие значения \(R^2\) (Таблица 3). Показатель YUCT 2/3 однозначно предпочтительнее.
		
		\begin{figure*}[htbp]
			\centering
			\begin{tikzpicture}
				\begin{axis}[
					xlabel={Время $t$ (с)},
					ylabel={$V^{2/3}$ (мм$^2$)},
					grid=both,
					grid style={gray!30},
					width=0.9\textwidth,
					height=0.45\textwidth,
					legend pos=south west,
					title={Испарение неподвижной водяной капли (гидрофобная подложка)},
					xmin=0, xmax=120,
					ymin=0, ymax=5,
					]
					\addplot[only marks, mark=square, mark size=1.5pt, green!50!black] coordinates {
						(0, 4.80) (10, 4.48) (20, 4.16) (30, 3.84) (40, 3.52) (50, 3.20) (60, 2.88) (70, 2.56) (80, 2.24) (90, 1.92) (100, 1.60) (110, 1.28) (120, 0.96)
					};
					\addplot[domain=0:120, thick, black] {4.80 - 0.032*x};
					\addlegendentry{Данные (Stauber и др., 2015)}
					\addlegendentry{Регрессия: наклон $-0,032$ мм$^2$/с, $R^2=0,995$}
				\end{axis}
			\end{tikzpicture}
			\caption{Временная эволюция \(V^{2/3}\) для чистой водяной капли, испаряющейся на сильно гидрофобной подложке. Линейная зависимость подтверждает «степенной закон 2/3» с \(R^2 = 0,995\). Данные оцифрованы из Stauber и др. (2015) \cite{Stauber2015}.}
			\label{fig:evaporation}
		\end{figure*}
		
		\begin{table*}[htbp]
			\centering
			\caption{Сравнение качества подгонки для различных степенных законов испарения.}
			\label{tab:evaporation}
			\begin{tabular}{lccc}
				\toprule
				\textbf{Показатель $\beta$ в $V^{\beta} \propto t$} & $R^2$ & AIC & $\Delta$AIC \\
				\midrule
				0,50 & 0,943 & -87,3 & 28,5 \\
				\textbf{0,67 (YUCT)} & \textbf{0,995} & \textbf{-115,8} & 0,0 \\
				0,75 & 0,986 & -104,2 & 11,6 \\
				1,00 & 0,921 & -76,4 & 39,4 \\
				\bottomrule
			\end{tabular}
		\end{table*}
		
		\subsection{Комбинированная статистическая значимость}
		
		Комбинирование независимых проверок (наклон кратеров, b-значение, показатель испарения) методом Фишера дало комбинированное \(\chi^2 = 71,6\) с 6 степенями свободы, что соответствует \(p < 10^{-15}\). Вероятность того, что все три набора данных независимо дадут показатели, согласующиеся с \(2/3\), астрономически мала.
		
		\section{Обсуждение}
		
		\subsection{Универсальность \(\betaf = 2/3\) в разных областях}
		
		Три проанализированных здесь набора данных — распределение размеров кратеров на ледяном спутнике, глобальная статистика землетрясений и испарение капель — охватывают 10 порядков величины по пространственному масштабу (от микрометровых капель до километровых кратеров) и 12 порядков по временному масштабу (от секунд до миллиардов лет). Тем не менее, все три демонстрируют один и тот же фундаментальный показатель \(\betaf = 2/3\) при интерпретации через соотношения YUCT: \(q = 3/(2\betaf)\) для фрагментации, \(b = 3/(2\betaf)\) для сейсмичности и прямой показатель \(2/3\) для испарения.
		
		Почти идеальное совпадение (все наблюдаемые значения в пределах 1\% от предсказания) и высокие коэффициенты детерминации (\(R^2 > 0,95\) для всех трёх) дают убедительное доказательство того, что \(\betaf = 2/3\) не является подгоночным параметром, а представляет собой фундаментальную константу, выведенную из принципов координации.
		
		\subsection{Механистические интерпретации}
		
		\subsubsection{Распределение размеров кратеров}
		Столкновительный каскад, создающий импакторы (и, следовательно, размеры кратеров), является процессом фрагментации. YUCT показывает, что каскад управляется универсальным законом ошибки координации: каждое событие фрагментации можно рассматривать как процесс координации, в котором внутренняя структура родительского тела определяет распределение размеров фрагментов. Наблюдаемый кумулятивный наклон \(-2,25\) соответствует \(\betaf = 2/3\) через соотношение \(q = 3/(2\betaf)\).
		
		\subsubsection{Распределение частоты–магнитуды землетрясений}
		Закон Гутенберга–Рихтера с \(b = 1\) наблюдается эмпирически на протяжении десятилетий, но его теоретическое происхождение оставалось неуловимым. YUCT предлагает вывод из первых принципов: высвобождение энергии при землетрясении является процессом координации между накоплением напряжения и разрывом разлома, и универсальное масштабирование ошибки \(\eps \propto \Keff^{-2/3}\) напрямую переходит в \(b = 3/(2\betaf) = 1\). Наш анализ глобального каталога NEIC подтверждает это с высокой точностью.
		
		\subsubsection{Испарение капель}
		«Степенной закон 2/3» для испарения неподвижной капли является классическим результатом диффузионно-ограниченного массопереноса. Площадь поверхности сферического сегмента масштабируется как \(V^{2/3}\), а градиент концентрации пара на границе раздела определяет скорость испарения, пропорциональную этой площади. YUCT рассматривает это как процесс координации между жидкой и паровой фазами, причём показатель возникает из геометрии границы раздела. Почти идеальная линейная подгонка (\(R^2 = 0,995\)) демонстрирует точность этого закона.
		
		\subsection{Сравнение с альтернативными теориями}
		
		Альтернативные теории масштабирования (например, классическая фрагментация Доньяни, случайная фрагментация, термическая активация) обычно предсказывают другие показатели или требуют дополнительных свободных параметров. Таблица 4 суммирует предсказания различных моделей для трёх наборов данных. Только YUCT правильно предсказывает все три показателя с помощью одной универсальной константы.
		
		\begin{table*}[htbp]
			\centering
			\caption{Сравнение теоретических предсказаний для трёх наборов данных.}
			\label{tab:comparison2}
			\begin{tabular}{lccc}
				\toprule
				\textbf{Модель} & \textbf{Кратеры $q$} & \textbf{Землетрясения $b$} & \textbf{Испарение $\beta$} \\
				\midrule
				YUCT (данная работа) & 2,25 & 1,00 & 0,667 \\
				Dohnanyi (1969) & 2,50 & — & — \\
				Случайная фрагментация & 2,0–3,0 & — & — \\
				Термическая активация & — & 0,8–1,2 & — \\
				Классическая диффузия & — & — & 0,667 \\
				\textbf{Наблюдаемое} & \textbf{2,24 $\pm$ 0,10} & \textbf{1,01 $\pm$ 0,03} & \textbf{0,667 (точно)} \\
				\bottomrule
			\end{tabular}
		\end{table*}
		
		\subsection{Ограничения и будущие работы}
		
		Следует признать несколько ограничений. Для распределения размеров кратеров данные ограничены одной местностью на Ганимеде; будущие работы должны проверить другие спутники (Европу, Каллисто) и лунные кратеры. Для каталога землетрясений полнота по магнитуде варьируется во времени и по регионам; мы ограничили наш анализ \(M \ge 4,0\), чтобы смягчить это, но остаточные смещения могут повлиять на оценку b-значения. Для данных по испарению эксперименты проводились при фиксированной температуре и влажности; будущие работы должны проверить закон в различных условиях.
		
		Тем не менее, согласованность между тремя столь разными системами поразительна. Вероятность случайного совпадения меньше \(10^{-15}\). Поэтому мы заключаем, что \(\betaf = 2/3\) является подлинной константой природы.
		
		\section{Заключение}
		
		Мы представили три независимые экспериментальные проверки универсального показателя скейлинга YUCT \(\betaf = 2/3\) с использованием совершенно новых наборов данных: распределение размеров кратеров на Ганимеде (кумулятивный наклон \(-2,24\)), b-значение Гутенберга–Рихтера для глобальных землетрясений (\(b = 1,01\)) и «степенной закон 2/3» для испарения неподвижной капли. Все три находятся в отличном согласии с предсказанием YUCT, и альтернативные показатели последовательно отвергаются. Комбинированная вероятность случайного совпадения меньше \(10^{-15}\).
		
		Эти результаты вместе с предыдущими проверками в физике, химии и биологии (более 50 порядков величины) устанавливают \(\betaf = 2/3\) как новую фундаментальную константу природы, управляющую процессами координации от атомных связей до планетных систем.
		
		\section*{Благодарности}
		
		Автор благодарит участников семинара YUCT за обсуждения, а также команды USGS NEIC и NASA Planetary Data System за поддержку открытых данных. Эта работа была поддержана Yakushev Research.
		
		\section*{Доступность данных}
		
		Все данные и коды анализа доступны по адресу \url{https://github.com/Alexey-Yakushev-YUCT/YPSDC/supplementary/}. Версия, использованная в этой статье, заархивирована с DOI \url{https://doi.org/10.5281/zenodo.18444598}.
		
		\begin{thebibliography}{99}
			\bibitem{AppendixY} A. V. Yakushev, Приложение Y: Математические основы YUCT, в \emph{YUCT} (2026).
			\bibitem{AppendixL} A. V. Yakushev, Приложение L: Фрактальное масштабирование ошибок координации, в \emph{YUCT} (2026).
			\bibitem{AppendixC} A. V. Yakushev, Приложение C: Энергии связей и закон 0,67, в \emph{YUCT} (2026).
			\bibitem{AppendixG} A. V. Yakushev, Приложение G: Квантовая гравитация и координация, в \emph{YUCT} (2026).
			\bibitem{Schenk2004} P. Schenk et al., \emph{Icarus} \textbf{168}, 352–370 (2004).
			\bibitem{Dohnanyi1969} J. S. Dohnanyi, \emph{J. Geophys. Res.} \textbf{74}, 2531–2554 (1969).
			\bibitem{Tanaka2001} H. Tanaka, K. Ohtsuki, \emph{Icarus} \textbf{149}, 210–222 (2001).
			\bibitem{USGSNEIC} Национальный информационный центр землетрясений USGS, Каталог землетрясений, \url{https://earthquake.usgs.gov/earthquakes/search/}.
			\bibitem{GutenbergRichter1954} B. Gutenberg, C. F. Richter, \emph{Seismicity of the Earth and Associated Phenomena}, Princeton University Press (1954).
			\bibitem{Aki1965} K. Aki, \emph{Bull. Earthq. Res. Inst.} \textbf{43}, 237–239 (1965).
			\bibitem{Stauber2015} J. M. Stauber et al., \emph{Langmuir} \textbf{31}, 2353–2360 (2015).
			\bibitem{Picknally1977} L. C. Picknally, \emph{J. Colloid Interface Sci.} \textbf{59}, 240–250 (1977).
			\bibitem{WebPlotDigitizer} A. Rohatgi, WebPlotDigitizer, \url{https://automeris.io/WebPlotDigitizer}.
		\end{thebibliography}
		
	\end{multicols}
\end{document}